% from ega4-8-fr.tex
\oldpage[IV-3]{48}
\begin{corollary}[8.12.9]
\label{IV.8.12.9-fr}
Soit $f:X\to Y$ un morphisme localement quasi-fini
$(\mathbf{Err}_{\mathrm{III}}, 20)$. Pour tout $x\in X$ il existe un
voisinage ouvert $U$ de $x$ dans $X$, un voisinage ouvert $V$ de
$y=f(x)$ dans $Y$, tels que $f(U)\subset V$ et une factorisation
\[
U\xrightarrow{f'}V'\xrightarrow{u}V
\]
de la restriction de $f$ à $U$, où $f'$ est une immersion ouverte et $u$ un
morphisme fini.
\end{corollary}

Il suffit évidemment de prendre pour $V$ un voisinage affine de $y$ dans
$Y$, pour $U$ un voisinage affine de $x$ dans $X$ contenu dans $f^{-1}(U)$
et tel que $f|U$ soit quasi-fini. Le morphisme $U\to V$ restriction de $f$
étant alors affine (donc quasi-projectif), on peut lui appliquer (8.12.8).

\begin{corollary}[8.12.10]
\EGAReaderTarget{ega4.c04.cor.8-12-10}{IV.8.12.10-fr}
Soient $Y$ un préschéma intègre et normal, $X$ un préschéma intègre,
$f:X\to Y$ un morphisme birationnel et localement quasi-fini
$(\mathbf{Err}_{\mathrm{III}}, 20)$. Alors $f$ est un isomorphisme local~;
pour que $f$ soit une immersion ouverte, il faut et il suffit que $f$ soit
en outre séparé.
\end{corollary}

La seconde assertion résulte aussitôt de la première et de (I, 8.2.8). Pour
prouver la première assertion, on peut supposer $X$ et $Y$ affines et $f$
quasi-fini~; considérons la factorisation $f=u\circ f'$ de (8.12.8), qui
permet d'identifier $X$ par $f'$ à un sous-préschéma induit sur un ouvert de
$Y'$. Comme $X$ est intègre, on peut, en vertu de (I, 5.2.3), remplacer
$Y'$ par le sous-préschéma réduit de $Y'$ ayant pour espace sous-jacent
$\overline X$, donc on peut supposer aussi que $Y'$ est \emph{intègre}. En
outre, puisque $f$ est birationnel, $u$ l'est aussi. La conclusion résulte
alors du lemme suivant~:

\begin{lemma}[8.12.10.1]
\label{IV.8.12.10.1-fr}
Soient $Y'$ un préschéma intègre, $Y$ un préschéma intègre et normal~; alors
un morphisme $u:Y'\to Y$ fini et birationnel est un isomorphisme.
\end{lemma}

Posons en effet $\mathcal A=u_*(\mathcal O_{Y'})$, de sorte que $\mathcal A$
est une $\mathcal O_Y$-Algèbre finie, $Y'$ s'identifiant à
$\operatorname{Spec}(\mathcal A)$ (II, 1.3.6). Si $R(Y)$ est le corps des
fonctions rationnelles de $Y$, on a donc, pour tout $y\in Y$,
$\mathcal O_{Y,y}\subset\mathcal A_y\subset R(Y)$~; mais comme l'anneau
$\mathcal O_{Y,y}$ est par hypothèse intégralement clos et a $R(Y)$ pour
corps des fractions, on a nécessairement
$\mathcal A_y=\mathcal O_{Y,y}$, d'où le lemme.

\begin{corollary}[8.12.11]
\EGAReaderTarget{ega4.c04.cor.8-12-11}{IV.8.12.11-fr}
Soient $Y$ un préschéma intègre, $X$ un préschéma intègre et normal,
$f:X\to Y$ un morphisme dominant et localement quasi-fini. Soient $K$ et
$L$ (extension de $K$) les corps de fonctions rationnelles de $Y$ et $X$
respectivement~; et soit $Y'$ la fermeture intégrale de $Y$ relativement à
$L$ (II, 6.3.4)~; alors $f$ se factorise d'une seule manière en
$f:X\xrightarrow{f'}Y'\xrightarrow{u}Y$, où $f'$ est birationnel, et
correspond à l'automorphisme identique de $L$~; $f'$ est alors un
isomorphisme local, et pour que $f'$ soit une immersion ouverte, il faut et
il suffit que $f$ soit séparé.
\end{corollary}

L'existence et l'unicité de la factorisation de $f$ résultent de (II, 6.3.9).
Il résulte de (II, 6.2.4, (v)), en se ramenant au cas affine, que $f'$ est
localement quasi-fini~; en outre, il résulte de (I, 5.5.1) que, pour que
$f'$ soit séparé, il faut et il suffit que $f$ le soit, puisque $u$ est
affine, donc séparé~; les deux dernières assertions sont donc des
conséquences de (8.12.10) appliqué à $f'$.

